\section{Physical validation plan}
\label{sec:validation}

No part of the hardware has been measured. Each row is a physical test still to
run. Pass criteria are the requirement thresholds in \ev{management/REQUIREMENTS.md};
where a requirement gives no threshold, the analysis prediction is the criterion
and is labelled. Items P-1 to P-6 are the systems plan
\ev{management/handoffs/systems-engineer/T-011.md}; hardware items from firmware are
in \ev{reports/firmware/coverage.md}; electronics next tests are in the
T-050 to T-055, T-073 and T-076 handoffs. All rows are NOT RUN.

{\footnotesize
\begin{longtable}{@{}>{\raggedright\arraybackslash}p{0.07\linewidth}>{\raggedright\arraybackslash}p{0.21\linewidth}>{\raggedright\arraybackslash}p{0.27\linewidth}>{\raggedright\arraybackslash}p{0.29\linewidth}>{\raggedright\arraybackslash}p{0.10\linewidth}@{}}
\caption{Physical validation plan (all NOT RUN).}\label{tab:validation}\\
\toprule
ID & What to measure & How / instrument & Pass criterion & REQ \\
\midrule
\endfirsthead
\toprule
ID & What to measure & How / instrument & Pass criterion & REQ \\
\midrule
\endhead
\bottomrule
\endlastfoot
V-01 (P-1) & Object plus lead capacitance & LCR meter on typical objects (bare pad, fruit, glass of water) with 20--50\,cm leads & Touch detected after recalibration from about 20\,pF to about 500\,pF (\lbl{ASSUMPTION} range in REQ-037) & REQ-037 \\
V-02 (P-2) & Touch-to-USB latency and jitter & Logic analyser on a test electrode, SDA/SCL and USB & Worst case $\le$ 15\,ms; jitter $\le$ 10\,ms, for $C \le 100$\,pF & REQ-014, 015 \\
V-03 (P-3) & VBUS current per Pico variant & USB inline current meter; Pico, Pico W, Pico 2, Pico 2 W & $\le$ 100\,mA & REQ-019 \\
V-04 & VBUS range & Adjustable bench supply at 4.40\,V and 5.50\,V at the connector & Device enumerates and plays at both ends & REQ-018 \\
V-05 (P-4) & I\textsuperscript{2}C rise time and low levels, PCB and breadboard & Oscilloscope on SDA/SCL at the Pico and at the module & Rise $\le$ 300\,ns at 400\,kHz; low levels below VIL (0.93\,V at the MPR121, 0.8\,V at the Pico) & REQ-026, 041 \\
V-06 & Misplaced data wire current & Wire J1 SDA/SCL to a supply hole; current meter in series & $\le$ 15.5\,mA per line (\lbl{CALCULATION}); within RP2040 IO limits & REQ-024 \\
V-07 (P-5) & Host enumeration & Windows 10, Windows 11, macOS 12 or later; MIDI monitor & No vendor driver; name ``FrutaMIDI''; Note On/Off and poly aftertouch on channel 1 only & REQ-011, 012, 013 \\
V-08 (P-6) & False triggers and crosstalk & 10\,min idle with all 12 electrodes on objects; then touch each pad and watch neighbours & No false Note On; no crosstalk & REQ-038 \\
V-09 & Electrode readings with and without $R$ & \texttt{SERIAL\_DEBUG} deltas per object, with the 3.3\,k$\Omega$ resistor and with it bridged & Touch detected on every object; delta ratio $\ge$ 0.904 (\lbl{SIMULATION} prediction) & REQ-037 \\
V-10 & ESD robustness & IEC 61000-4-2 contact discharge $\pm$4\,kV on each terminal, at a test lab & No permanent damage: MPR121 still responds, all pads still detect touch, baseline drift $<$ 1 count (T-060) & REQ-023 \\
V-11 & Series resistor after ESD & Watch for flashover across R1--R12 during V-10; measure R1--R12 afterwards & No flashover; each resistance within 1\,\% of its pre-test value (T-060) & REQ-023 \\
V-12 & JP1 bridge and etch gaps & Visual inspection with a magnifier and continuity after etching and bridging & Bridge only on the chosen side; the 1.87\,mm gap on the other side clear (ADR-021) and 0.6\,mm track gaps clear; no short to GND in any position & REQ-025, 039 \\
V-13 & Real loop period & Toggle a spare GPIO each loop; oscilloscope & No blocking $>$ 10\,ms during play, except user recalibration and the lost-sensor restart (ADR-018) & REQ-016 \\
V-14 & Recalibration duration & GPIO toggle around \texttt{recalibrate()} & Recorded per sensor; compare with 16.8--22.1\,ms (\lbl{CALCULATION}); no stuck notes afterwards & REQ-008, 002 \\
V-15 & Pot ADC noise on the PCB & Log \texttt{analogRead} at each zone border for 10\,s & Peak-to-peak noise $<$ 41 counts & REQ-006 \\
V-16 & Blink with live pads & Touch pads while the LED blinks scale 6; press again during a blink & Notes sent during the blink; pattern restarts with the new scale & REQ-007, 016 \\
V-17 & Button bounce and long press & Oscilloscope on the button pin; press tests & Bounce $<$ 10\,ms (\lbl{ASSUMPTION}); long press acts after $>$ 800\,ms & REQ-007, 008 \\
V-18 & Sensor absent or lost, and LED & Power without the module; then with it; unplug the module during play; every Pico board & No MIDI from the missing sensor, its notes released, LED at 5\,Hz, retry every 2\,s; 2 slow blinks when found; LED lit while touched (W boards via CYW43) & REQ-009, 010 \\
V-19 & MPR121 configuration and burst reads & Read back 0x7B--0x7F after start-up; logic analyser during play & Values as configured; 2 transactions per sensor per aftertouch pass; burst data equal to per-electrode reads & REQ-004, 037 \\
V-20 & Sampler in real browsers & Chrome and Edge on Windows and macOS with the device & Pads 36--47 sound; velocity 0 acts as Note Off; learn swaps notes & REQ-034 \\
V-21 & Sampler with a busy device & Open the device in a DAW on Windows first, then the sampler & Sampler reports the device as busy & REQ-034 \\
V-22 & Shared notes & Set an octave so two pads share a note; touch both, release one, then the other & One Note On; Note Off only after the last release; aftertouch = the stronger pad & REQ-042 \\
V-23 & v1.2 configuration wires (R-008) & Inspect the two J1 $\rightarrow$ S wires before power; continuity J1 SDA/SCL to the S hole of the module revision & Each wire on the correct S hole (S5/S6 original, S3/S4 STEMMA QT), flat and mechanically secure & REQ-024, 025 \\
V-24 & R1--R12 fitted as designed & Inspect the resistor row after soldering; caliper on one resistor & Lying flat, part MBE04140C3301FC100, leads 15.24\,mm apart, no lead bent towards a neighbour; leads formed on a gauge, the two straight lengths equal within about 0.25\,mm (T-073) & REQ-023 \\
V-25 & Arc path and I\textsuperscript{2}C upsets during ESD (C2-005) & During V-10, watch each terminal for an arc to the GND edge track; after every strike check for I\textsuperscript{2}C NACKs (logic analyser on SDA/SCL), the lost-sensor LED blink, Pico resets (USB re-enumeration) and stuck notes (MIDI monitor) & Arc, if any, to the GND track only (not across R$x$); no reset, no lost sensor and no stuck note after any strike; any I\textsuperscript{2}C error recorded and reported (ADR-020) & REQ-023 \\
V-26 & MBE 0414 body and end-cap length (C2-002) & Caliper on one resistor before fitting: body length $L$ and cap length $L_c$; then re-run \nolinkurl{check_gap.py} section (d) with the measured values & $L_c \le$ 3.0\,mm (4\,kV rule met at every swept $L$ and placement error, \lbl{CALCULATION}); otherwise the recomputed cap gap $\ge$ 3.5\,mm at the measured placement error & REQ-023 \\
V-27 & I\textsuperscript{2}C traffic against electrode baseline noise (C2-012) & \texttt{SERIAL\_DEBUG} deltas per electrode with the bus idle between reads versus continuous I\textsuperscript{2}C polling, electrodes on typical objects & No false Note On; noise with traffic increase of at most 3 counts in the per-electrode delta (\lbl{ASSUMPTION} threshold, T-074) & REQ-038 \\
V-28 & ADR-021 parts and pads & Test-fit C1 (5.08\,mm leads) in its 0.8\,mm holes; bridge JP1 with an R13--R16 lead off-cut over the 1.87\,mm gap; tighten and re-tighten every J2 screw; with the pot fitted, log \texttt{analogRead} during slow and fast turns & C1 seated flat without force; JP1 bridge on one side only, no continuity to the other side; no J2 pad lifted or cracked; the octave follows the knob with no visible lag (settling $\le$ 2.89\,ms, \lbl{CALCULATION}) & REQ-006, 025, 029 \\
\end{longtable}
}
